\documentclass[11pt]{article}

\usepackage[margin=1in]{geometry}
\usepackage{amsmath,amssymb,amsthm,mathtools}
\usepackage{microtype}
\usepackage[hidelinks]{hyperref}
\usepackage{enumitem}
\usepackage{xcolor}

\definecolor{deepgreen}{HTML}{17352E}
\definecolor{softgreen}{HTML}{EAF3EE}
\definecolor{warmgray}{HTML}{5B6663}

\newtheorem{proposition}{Proposition}
\newtheorem{lemma}[proposition]{Lemma}
\theoremstyle{remark}
\newtheorem{remark}[proposition]{Remark}

\newcommand{\E}{E(2)}
\newcommand{\Z}{S^2\times S^2}

\hypersetup{
  pdftitle={A Scope Observation Concerning K3 Problem 4.8},
  pdfauthor={},
  pdfsubject={Fintushel--Stern knot surgery and stabilization}
}

\title{A Scope Observation Concerning K3 Problem 4.8\\[4pt]
\large The displayed arbitrary-$X$ formulation is already negative after one stabilization}
\author{}
\date{}

\begin{document}
\maketitle
\vspace{-1.5em}

\begin{center}
\fcolorbox{deepgreen}{softgreen}{\parbox{0.89\textwidth}{
\textbf{Purpose of this note.}
This note reports a possible scope issue in Problem 4.8 of
\emph{K3: A New Problem List in Low-Dimensional Topology}.
The underlying stabilization theorem is not new.  The observation is that
its direct application appears to give a negative answer to the problem as
displayed for arbitrary qualifying $X$.  This does not answer the separate
Fintushel--Stern question with $X$ fixed to the K3 surface $E(2)$.
}}
\end{center}

\section{The displayed problem and the K3 special case}

Problem 4.8 fixes a closed simply connected smooth $4$--manifold $X$ and a
smoothly embedded torus $T\subset X$ satisfying
\[
  \pi_1(X-T)=1,
  \qquad [T]^2=0.
\]
For a knot $K\subset S^3$, let $X_K$ denote Fintushel--Stern knot surgery
along $T$.  The displayed question asks whether, for prime knots $K_1,K_2$,
a diffeomorphism $X_{K_1}\cong X_{K_2}$ forces $K_1$ to be $K_2$ or its
mirror.

As written, $X$ is arbitrary subject to the stated hypotheses.  In
particular, the statement imposes no minimality, irreducibility, or
unstabilized condition.  Remark (2) then identifies $X=K3=E(2)$ as a
particular case, attributes that case to Fintushel and Stern, and states that
beyond equality of Alexander polynomials no further constraint is known.
The distinction between the displayed arbitrary-$X$ question and this fixed
K3 case is essential.

\section{The previously published stabilization theorem}

Let $Y$ contain a c-embedded torus, so that the knot-surgery construction is
performed in a cusp neighborhood.  Akbulut's Theorem 2.2 states that for
every knot $K$,
\begin{equation}\label{eq:akbulut}
  Y_K\mathbin{\#}(S^2\times S^2)
  \cong
  Y\mathbin{\#}(S^2\times S^2).
\end{equation}
The paper notes that Auckly independently observed the same stabilization
phenomenon.  Baykur later proved a one-stabilization theorem under the more
general hypotheses of a square-zero torus with simply connected complement.

Equation~\eqref{eq:akbulut} is an established theorem.  No new
stabilization result is claimed here.

\section{The direct consequence for the displayed formulation}

\begin{proposition}\label{prop:scope}
The displayed arbitrary-$X$ formulation of K3 Problem 4.8 has a negative
answer.
\end{proposition}

\begin{proof}
Let $Y=E(2)$ and let $F\subset Y$ be a regular elliptic fiber.  The fiber is
a c-embedded square-zero torus, and its complement is simply connected.
Put
\[
  Z=S^2\times S^2,
  \qquad X=Y\# Z,
\]
where the connected sum is performed in a ball disjoint from $F$.
Then $X$ is closed, simply connected, and smooth; the same copy of $F$ has
$F^2=0$ and $\pi_1(X-F)=1$.

Knot surgery is disjoint from the connected-sum neck.  Regrouping the same
cut-and-paste pieces gives an orientation-preserving locality
diffeomorphism
\begin{equation}\label{eq:locality}
  X_K\cong Y_K\# Z
\end{equation}
for every knot $K$.  Applying Akbulut's theorem to $(Y,F)$ gives
\begin{equation}\label{eq:stable}
  Y_K\# Z\cong Y\# Z=X.
\end{equation}
Combining \eqref{eq:locality} and \eqref{eq:stable},
\[
  X_K\cong X
\]
for every knot $K$.

Now choose
\[
  K_1=T(2,3),
  \qquad K_2=T(2,5).
\]
Both are prime torus knots.  Their Alexander polynomials have respective
breadths $2$ and $4$, so they are neither equivalent nor mirrors.  Yet
\[
  X_{K_1}\cong X\cong X_{K_2}.
\]
This is a counterexample satisfying every hypothesis in the displayed
arbitrary-$X$ formulation.
\end{proof}

\section{What the observation does and does not establish}

\begin{itemize}[leftmargin=2em,itemsep=0.45em]
  \item It gives a direct negative answer to Problem 4.8 exactly as displayed
  when arbitrary qualifying ambient manifolds are permitted.
  \item It does not solve the unstabilized fixed-$E(2)$ question in
  Remark (2).
  \item It does not introduce a new $4$--manifold theorem: the decisive
  input is Akbulut's published Theorem 2.2, or alternatively Baykur's later
  one-stabilization theorem.
  \item The potentially new point is the explicit application of that old
  theorem to the scope of the newly published Problem 4.8.  A targeted
  search found no public correction or prior explicit notice of this
  implication, but no claim of worldwide priority is made without editorial
  confirmation.
\end{itemize}

\section{Suggested editorial clarification}

If the intended problem is the Fintushel--Stern question for the K3 surface,
the most direct repair would be to state the ambient pair explicitly, for
example:

\begin{quote}
Let $X=E(2)$ and let $T$ be a regular elliptic fiber.  If $K_1$ and $K_2$
are prime knots such that $X_{K_1}$ and $X_{K_2}$ are diffeomorphic, must
$K_1$ be $K_2$ or its mirror?
\end{quote}

More general unstabilized variants may also be meaningful, but additional
hypotheses would need to be selected by the authors; minimality or
irreducibility should not be inserted without a separate analysis.

The editors and scribes are respectfully asked to clarify:
\begin{enumerate}[leftmargin=2em,itemsep=0.35em]
  \item whether the arbitrary-$X$ formulation was intended;
  \item whether the implication of Akbulut's Theorem 2.2 was already known
  in connection with Problem 4.8; and
  \item whether a correction restricting the problem to $X=E(2)$, or another
  explicitly unstabilized class, is appropriate.
\end{enumerate}

\section{Verification materials}

An explicit certificate, source ledger, independent replay program, Lean 4
formalization of the logical composition, and hash-bound kernel evidence are
available at
\begin{center}
  \url{https://open.argusbot.cn/results/2884}.
\end{center}
The Lean development checks the composition from the ambient, locality, and
stabilization hypotheses to the counterexample witness.  It does not
formalize smooth $4$--manifold topology or Akbulut's theorem from first
principles; those inputs remain explicit hypotheses tied to the primary
sources.

\begin{thebibliography}{9}

\bibitem{Akbulut2002}
S. Akbulut,
\emph{Variations on Fintushel--Stern knot surgery on 4-manifolds},
Turkish J. Math. \textbf{26} (2002), 81--92.
\url{https://arxiv.org/abs/math/0201156}.

\bibitem{Baykur2018}
R. I. Baykur,
\emph{Dissolving knot surgered 4-manifolds by classical cobordism arguments},
J. Knot Theory Ramifications \textbf{27} (2018), no. 5, 1871001.
\url{https://arxiv.org/abs/1704.04491}.

\bibitem{K3}
R. I. Baykur, R. C. Kirby, and D. Ruberman,
\emph{K3: A New Problem List in Low-Dimensional Topology},
AMS Surveys and Monographs, vol. 295, 2026, Problem 4.8.
\url{https://doi.org/10.1090/surv/295}.

\bibitem{FS1998}
R. Fintushel and R. J. Stern,
\emph{Knots, links, and 4-manifolds},
Invent. Math. \textbf{134} (1998), 363--400.
\url{https://arxiv.org/abs/dg-ga/9612014}.

\end{thebibliography}

\end{document}
